\documentclass[11pt]{article}

\usepackage[a4paper,margin=1.05in]{geometry}
\usepackage[T1]{fontenc}
\usepackage[utf8]{inputenc}
\usepackage{lmodern}
\usepackage{microtype}
\usepackage{amsmath,amssymb,amsthm,mathtools}
\usepackage{booktabs}
\usepackage{enumitem}
\usepackage{xcolor}
\usepackage[colorlinks=true,linkcolor=blue!45!black,citecolor=blue!45!black,urlcolor=blue!45!black]{hyperref}

\numberwithin{equation}{section}

\newtheorem{theorem}{Theorem}[section]
\newtheorem{proposition}[theorem]{Proposition}
\newtheorem{lemma}[theorem]{Lemma}
\newtheorem{corollary}[theorem]{Corollary}
\newtheorem{definition}[theorem]{Definition}
\newtheorem{remark}[theorem]{Remark}

\newcommand{\C}{\mathbb C}
\newcommand{\N}{\mathbb N}
\newcommand{\Sym}{\mathfrak S}
\newcommand{\cyc}{\#}
\newcommand{\tr}{\operatorname{Tr}}
\newcommand{\id}{\mathrm{id}}
\newcommand{\g}{\gamma}
\newcommand{\dd}{\delta}
\newcommand{\ddp}{\delta_{+}}
\newcommand{\ddm}{\delta_{-}}
\newcommand{\rank}{\mathcal R}
\newcommand{\Count}{\mathcal C}
\newcommand{\Countp}{\mathcal C^{+}}
\newcommand{\Env}{\mathcal E}
\newcommand{\Loss}{Q}
\newcommand{\Wish}{W}
\newcommand{\PT}{\Gamma}

\title{Bidefect Counting for the Aubrun Partial-Transpose Moment Expansion}
\author{Kieran McShane}
\date{}

\begin{document}

\maketitle

\begin{abstract}
Aubrun's moment expansion for the partial transpose of a random induced
state is governed by two Cayley geodesic inequalities for the long cycle.
The natural bookkeeping variable is therefore not a single total defect, but
the ordered pair of defects associated with the two sides of the long cycle.
We give a self-contained pure-math account of this bidefect reduction.  The
rank loss is exactly the sum of the two oriented defects, the scalar sum is
absorbed by two shifted analytic bases, and the required count is reduced to a
one-sided genus enumeration for unicellular maps.  The latter is bounded by
the planar noncrossing count and Chapuy's slicing recurrence.  The result is a
clean combinatorial input to the balanced partial-transpose moment method once
the standard Wick expansion has been written in its two-base cycle-count form.
\end{abstract}

\tableofcontents

\section{Introduction}

Let \(\Wish\) be a complex Wishart matrix acting on
\(\C^d\otimes\C^d\), and let \(\Wish^\PT\) denote the partial transpose on
one tensor factor.  The moment method for \(\Wish^\PT\) expands
\(\mathbb E\tr((\Wish^\PT)^k)\) as a finite sum indexed by permutations.
In the balanced regime \(s\asymp d^2\), the exponent of each Wick term is
controlled by the cycle counts of a permutation and of its products with the
long cycle.  This is the combinatorial core of Aubrun's analysis of the
PPT threshold for random induced states
\cite{aubrun_partial_2012,aubrun_entanglement_2014}.

The point of this paper is a bookkeeping correction that becomes important
when one wants a sharp, reusable estimate.  There are two Cayley geodesic
inequalities, one involving the long cycle and one involving its inverse.
Collapsing them into a single total defect loses the two-base structure of
the analytic term.  Keeping them separate gives two defects, \(\ddp\) and
\(\ddm\).  The plus defect is paid by one shifted base, the minus defect by
the other.  The scalar sum then becomes an ordinary double binomial
absorption.

The reduction has three parts.

\begin{enumerate}[label=(\roman*),leftmargin=*]
  \item The rank identity:
  \[
    \rank(\pi)+\ddp(\pi)+\ddm(\pi)=2n+2.
  \]
  This is an exact consequence of the two Cayley inequalities.
  \item The scalar absorption:
  \[
    \sum_{a,b}\Count_m(a,b)d^{m+3-a}S^{m+1-b}
    \le
    4^{m+1}(d+\Loss_m)^{m+3}(S+\Loss_m)^{m+1}.
  \]
  \item The count estimate: the bidefect count follows from a one-sided
  plus-defect count, and the latter follows from the unicellular-map genus
  formula together with the planar count and Chapuy slicing.
\end{enumerate}

The statements below are written in the language of permutations, Cayley
length, noncrossing partitions, and unicellular maps.

We do not rederive the full Wishart Wick formula.  The matrix-model
translation is isolated in Section~\ref{sec:moment-input}: the theorem there
assumes the standard termwise two-base form and then inserts the bidefect
count proved in the preceding sections.

\section{Cayley Defects}

Fix \(n\ge1\), write \(\Sym_n\) for the symmetric group on
\(\{0,\ldots,n-1\}\), and let
\[
  \g=(0\,1\,\cdots\,n-1).
\]
For \(\sigma\in\Sym_n\), let \(\cyc(\sigma)\) be the number of cycles of
\(\sigma\), fixed points included, and let
\[
  |\sigma|=n-\cyc(\sigma)
\]
be the Cayley length with respect to transpositions.

\begin{definition}[Oriented Cayley defects]
For \(\pi\in\Sym_n\), define
\[
  \ddp(\pi)=n+1-\cyc(\pi)-\cyc(\g\pi),
  \qquad
  \ddm(\pi)=n+1-\cyc(\pi)-\cyc(\pi\g^{-1}).
\]
\end{definition}

\begin{lemma}[Defects are nonnegative]\label{lem:defects-nonnegative}
For every \(\pi\in\Sym_n\),
\[
  \ddp(\pi)\ge0,
  \qquad
  \ddm(\pi)\ge0.
\]
\end{lemma}

\begin{proof}
The triangle inequality in the Cayley graph gives
\[
  |\pi|+|\pi^{-1}\g^{-1}|\ge |\g^{-1}|=n-1.
\]
Since \(\cyc(\pi^{-1}\g^{-1})=\cyc(\g\pi)\), this is
\[
  n-\cyc(\pi)+n-\cyc(\g\pi)\ge n-1,
\]
which is equivalent to \(\ddp(\pi)\ge0\).  The proof for \(\ddm\) is the
same, using
\[
  |\pi|+|\pi^{-1}\g|\ge |\g|=n-1
\]
and the identities
\[
  \cyc(\pi^{-1}\g)=\cyc(\g^{-1}\pi)=\cyc(\pi\g^{-1}).
\]
\end{proof}

\begin{lemma}[Support]\label{lem:defect-support}
For every \(\pi\in\Sym_n\),
\[
  0\le \ddp(\pi)\le n-1,
  \qquad
  0\le \ddm(\pi)\le n-1.
\]
\end{lemma}

\begin{proof}
The lower bounds are Lemma~\ref{lem:defects-nonnegative}.  For the upper
bounds, every permutation has at least one cycle.  Thus
\[
  \cyc(\pi)\ge1,\qquad
  \cyc(\g\pi)\ge1,\qquad
  \cyc(\pi\g^{-1})\ge1,
\]
and therefore
\[
  \ddp(\pi)=n+1-\cyc(\pi)-\cyc(\g\pi)\le n-1.
\]
The proof for \(\ddm\) is identical.
\end{proof}

\begin{definition}[Aubrun rank exponent]
The permutation rank exponent is
\[
  \rank(\pi)=2\cyc(\pi)+\cyc(\g\pi)+\cyc(\pi\g^{-1}).
\]
\end{definition}

\begin{proposition}[Rank identity]\label{prop:rank-identity}
For every \(\pi\in\Sym_n\),
\[
  \rank(\pi)+\ddp(\pi)+\ddm(\pi)=2n+2.
\]
Consequently \(\rank(\pi)\le 2n+2\).
\end{proposition}

\begin{proof}
Add the two definitions of the oriented defects:
\[
\begin{aligned}
  \ddp(\pi)+\ddm(\pi)
  &=
  2n+2
  -2\cyc(\pi)
  -\cyc(\g\pi)
  -\cyc(\pi\g^{-1})\\
  &=2n+2-\rank(\pi).
\end{aligned}
\]
The inequality follows from Lemma~\ref{lem:defects-nonnegative}.
\end{proof}

\section{Bidefect Slices and Scalar Absorption}

From now on set \(n=m+1\).  The bidefect slice is
\[
  \Count_m(a,b)
  =
  \#\{\pi\in\Sym_{m+1}:\ddp(\pi)=a,\ \ddm(\pi)=b\}.
\]
We use the polynomial loss
\[
  \Loss_m=(4(m+1))^{36}.
\]
The exponent \(36\) is deliberately generous: the genus argument below needs
only a fixed polynomial power, while the moment expansion around it may
require additional harmless polynomial losses.  Nothing in the scalar
argument depends on the precise value \(36\), only on using the same
\(\Loss_m\) consistently.  For the counting theorem in this paper, any fixed
power large enough to dominate \((2(m+1))^{3/2}\) per unit of plus defect
would suffice.

By Lemma~\ref{lem:defect-support}, the true support of the bidefect count is
contained in \(0\le a,b\le m\).  The slightly larger rectangle in
Definition~\ref{def:envelope} is the one used by the shifted-base binomial
absorption; it contains the full support.

\begin{remark}[Small cases]
Direct enumeration gives the following first bidefect slices:
\[
\begin{array}{ccl}
\toprule
m & n=m+1 & \text{nonzero values of }\Count_m(a,b)\\
\midrule
0 & 1 & \Count_0(0,0)=1\\
1 & 2 & \Count_1(0,0)=2\\
2 & 3 & \Count_2(0,0)=4,\quad
        \Count_2(2,0)=\Count_2(0,2)=1\\
3 & 4 & \Count_3(0,0)=9,\quad
        \Count_3(2,0)=\Count_3(0,2)=5,\quad
        \Count_3(2,2)=5\\
\bottomrule
\end{array}
\]
The table is not used below.  It records the first instances of the
two-defect support and makes visible why a single total defect is a coarser
bookkeeping variable.
\end{remark}

\begin{definition}[Bidefect envelope]\label{def:envelope}
For \(0\le a\le m+3\) and \(0\le b\le m+1\), set
\[
  \Env_m(a,b)
  =
  4^{m+1}
  \binom{m+3}{a}
  \binom{m+1}{b}
  \Loss_m^{a+b}.
\]
\end{definition}

\begin{theorem}[Scalar absorption]\label{thm:scalar-absorption}
Assume that, for all \(0\le a\le m+3\) and \(0\le b\le m+1\),
\[
  \Count_m(a,b)\le \Env_m(a,b).
\]
Then for all \(d,S\ge0\),
\[
  \sum_{a=0}^{m+3}\sum_{b=0}^{m+1}
    \Count_m(a,b)d^{m+3-a}S^{m+1-b}
  \le
  4^{m+1}(d+\Loss_m)^{m+3}(S+\Loss_m)^{m+1}.
\]
\end{theorem}

\begin{proof}
Apply the assumed bound term by term:
\[
\begin{aligned}
&\sum_{a=0}^{m+3}\sum_{b=0}^{m+1}
    \Count_m(a,b)d^{m+3-a}S^{m+1-b} \\
&\quad\le
4^{m+1}
\sum_{a=0}^{m+3}\binom{m+3}{a}\Loss_m^a d^{m+3-a}
\cdot
\sum_{b=0}^{m+1}\binom{m+1}{b}\Loss_m^b S^{m+1-b}.
\end{aligned}
\]
The displayed right side is a product of two binomial expansions, hence
\[
  4^{m+1}(d+\Loss_m)^{m+3}(S+\Loss_m)^{m+1}.
\]
\end{proof}

\begin{remark}
This is exactly why the two defects must not be collapsed.  A one-parameter
defect estimate naturally produces a common base.  The partial-transpose
moment term has two distinct analytic bases, and the bidefect count is the
combinatorial form compatible with those bases.
\end{remark}

\section{The Plus-Defect Genus}

The bidefect envelope follows from a one-sided count.  Define
\[
  \Countp_m(a)=\#\{\pi\in\Sym_{m+1}:\ddp(\pi)=a\}.
\]
Plainly \(\Count_m(a,b)\le \Countp_m(a)\).  We now bound
\(\Countp_m(a)\) using the unicellular-map interpretation of \(\ddp\).

\subsection{The map associated with a permutation}

Let \(n=m+1\).  We use the standard permutation model for rooted orientable
hypermaps: a triple
\[
  (\sigma_0,\sigma_1,\sigma_\infty)
\]
of permutations of \(\{0,\ldots,n-1\}\) with
\(\sigma_0\sigma_1\sigma_\infty=\id\) encodes a connected oriented hypermap
whenever the generated action is transitive.  The cycles of \(\sigma_0\) and
\(\sigma_1\) give the two vertex colors, the cycles of \(\sigma_\infty\) give
the faces, and the \(n\) labels are the darts.  The Euler characteristic is
\[
  \cyc(\sigma_0)+\cyc(\sigma_1)+\cyc(\sigma_\infty)-n=2-2g.
\]

To \(\pi\in\Sym_n\), associate the triple
\[
  \bigl(\pi^{-1},\,\g\pi,\,(\pi^{-1}\g\pi)^{-1}\bigr).
\]
The product of the three permutations is the identity.  The last permutation
is \(\pi^{-1}\g^{-1}\pi\), conjugate to \(\g^{-1}\), and hence has one
cycle.  Since this one face is an \(n\)-cycle on the dart labels, the group
generated by the three permutations acts transitively on the darts.  Thus the
associated hypermap is connected and unicellular.  The cycles of
\(\pi^{-1}\) give one vertex color and the cycles of \(\g\pi\) give the
other.

\begin{proposition}[Euler formula]\label{prop:euler}
The plus defect is twice the genus of this unicellular map:
\[
  \ddp(\pi)=2g(\pi).
\]
In particular \(\Countp_m(a)=0\) when \(a\) is odd.
\end{proposition}

\begin{proof}
The associated hypermap has \(n\) darts, one face, and
\[
  V=\cyc(\pi^{-1})+\cyc(\g\pi)=\cyc(\pi)+\cyc(\g\pi)
\]
vertices.  Euler's formula for an orientable unicellular map gives
\[
  V-n+1=2-2g(\pi).
\]
Therefore
\[
  2g(\pi)=n+1-\cyc(\pi)-\cyc(\g\pi)=\ddp(\pi).
\]
\end{proof}

\subsection{Planar count}

\begin{lemma}[Planar plus-defect count]\label{lem:planar}
For every \(m\ge0\),
\[
  \Countp_m(0)\le 4^{m+1}.
\]
\end{lemma}

\begin{proof}
The condition \(\ddp(\pi)=0\) is equality in the Cayley triangle inequality
between \(\id\), \(\pi\), and \(\g^{-1}\).  Equivalently, \(\pi\) belongs to
the Cayley interval \([\id,\g^{-1}]\).  By Biane's geodesic classification,
this interval is in bijection with the noncrossing partitions of the
\(n\)-gon \cite{biane_properties_1997}.  Its cardinality is the Catalan number
\[
  C_n=\frac{1}{n+1}\binom{2n}{n}.
\]
Since \(C_n\le 4^n=4^{m+1}\), the result follows.
\end{proof}

\subsection{Chapuy slicing}

Let
\[
  G_m(g)=\#\{\pi\in\Sym_{m+1}:g(\pi)=g\}.
\]
By Proposition~\ref{prop:euler}, \(G_m(g)=\Countp_m(2g)\).

\begin{lemma}[Chapuy slicing recurrence]\label{lem:slicing}
For every \(m\ge0\) and \(g\ge0\),
\[
  G_m(g+1)\le (2(m+1))^3G_m(g).
\]
\end{lemma}

\begin{proof}
Chapuy's trisection decomposition for unicellular maps says that slicing a
chosen trisection of a genus \(g+1\) unicellular map produces a genus \(g\)
unicellular map with three distinguished corners, and that the inverse gluing
operation recovers the original map with its chosen trisection
\cite{chapuy_new_2011}.  In the present labeled hypermap model, labels are carried
through the slicing and gluing operations, and unicellularity is preserved.
Root the face tour at dart \(0\) and order corners by first appearance along
the oriented boundary walk, breaking any remaining ties by the ordinary order
of labels.  Choose, for each positive-genus map, the first trisection in this
order.  Slicing that chosen trisection injects the objects counted by
\(G_m(g+1)\) into genus \(g\) objects with three marked corners: if two maps
had the same sliced map and the same three marked corners, the inverse gluing
would recover the same original map and the same chosen trisection.

The unicellular boundary tour uses \(2n\) oriented dart-sides, so there are at
most \(2n\) possible corners in each marked position.  Hence the target set
has cardinality at most \((2n)^3G_m(g)\), proving the recurrence.
\end{proof}

\begin{theorem}[Plus-defect envelope]\label{thm:plus-envelope}
For every \(m\ge0\) and \(a\ge0\),
\[
  \Countp_m(a)\le 4^{m+1}\Loss_m^a.
\]
\end{theorem}

\begin{proof}
If \(a\) is odd, the left side is zero by Proposition~\ref{prop:euler}.
Write \(a=2g\).  Iterating Lemma~\ref{lem:slicing} and using
Lemma~\ref{lem:planar} gives
\[
  \Countp_m(2g)=G_m(g)
  \le
  (2(m+1))^{3g}G_m(0)
  \le
  4^{m+1}(2(m+1))^{3g}.
\]
Since \(g=a/2\) and \(\Loss_m=(4(m+1))^{36}\ge(2(m+1))^{3/2}\), we have
\[
  (2(m+1))^{3g}=(2(m+1))^{3a/2}\le \Loss_m^a.
\]
\end{proof}

\begin{corollary}[Bidefect count]\label{cor:bidefect-count}
For \(0\le a\le m+3\) and \(0\le b\le m+1\),
\[
  \Count_m(a,b)\le \Env_m(a,b).
\]
\end{corollary}

\begin{proof}
Forgetting the minus defect gives
\[
  \Count_m(a,b)\le \Countp_m(a).
\]
Theorem~\ref{thm:plus-envelope} yields
\[
  \Count_m(a,b)\le 4^{m+1}\Loss_m^a.
\]
Since \(\binom{m+3}{a}\ge1\), \(\binom{m+1}{b}\ge1\), and
\(\Loss_m^b\ge1\), the right side is at most \(\Env_m(a,b)\).
\end{proof}

\begin{theorem}[Bidefect scalar bound]\label{thm:main}
For all \(m\ge0\) and all \(d,S\ge0\),
\[
  \sum_{a=0}^{m+3}\sum_{b=0}^{m+1}
    \Count_m(a,b)d^{m+3-a}S^{m+1-b}
  \le
  4^{m+1}(d+\Loss_m)^{m+3}(S+\Loss_m)^{m+1}.
\]
\end{theorem}

\begin{proof}
Combine Corollary~\ref{cor:bidefect-count} with
Theorem~\ref{thm:scalar-absorption}.
\end{proof}

\section{Consequence for the Partial-Transpose Moment Expansion}
\label{sec:moment-input}

We now state precisely what Theorem~\ref{thm:main} supplies to the
random-matrix estimate.  The precise normalization of a Wishart matrix varies
across the literature; the statement therefore abstracts the termwise
cycle-count input and does not claim to rederive the full Wick formula.

\begin{definition}[Two-base Wick form]
Fix \(k=m+1\) and set \(S=\sqrt{s}\).  A partial-transpose moment expansion
has the two-base bidefect form at length \(k\) if, after extracting the
normalization factor independent of \(\pi\), the magnitude of the term indexed
by \(\pi\in\Sym_k\) is bounded by
\[
  d^{m+3-\ddp(\pi)}S^{m+1-\ddm(\pi)}.
\]
Equivalently, the two cycle-count combinations appearing in the Wick exponent
are
\[
  \cyc(\pi)+\cyc(\g\pi)=k+1-\ddp(\pi),
  \qquad
  \cyc(\pi)+\cyc(\pi\g^{-1})=k+1-\ddm(\pi),
\]
with the model-dependent fixed offsets already absorbed into the exponents
\(m+3\) and \(m+1\).
\end{definition}

\begin{remark}
The definition is an explicit hypothesis on the chosen normalization of the
Wishart expansion.  In the usual Aubrun partial-transpose cycle-count formula
the \(\pi\)-dependent parts are precisely the two combinations displayed
above; trace normalization and the choice of \(S=\sqrt{s}\) account for the
fixed shifts in the powers \(m+3\) and \(m+1\)
\cite{aubrun_partial_2012,aubrun_entanglement_2014}.  Before using
Theorem~\ref{thm:aubrun-input} in another normalization, these fixed offsets
must be checked directly from the corresponding Wick formula.
\end{remark}

\begin{theorem}[Combinatorial input under the two-base Wick form]\label{thm:aubrun-input}
Let \(k=m+1\).  Suppose the length-\(k\) partial-transpose moment expansion
has the two-base bidefect form.  Then the full permutation sum, after the
normalization factor independent of \(\pi\) has been extracted, is bounded by
\[
  4^k(d+\Loss_{k-1})^{k+2}(S+\Loss_{k-1})^k,
  \qquad S=\sqrt{s}.
\]
Here the loss is polynomial in \(k\):
\[
  \Loss_{k-1}=(4k)^{36}.
\]
\end{theorem}

\begin{proof}
Group the terms by \(a=\ddp(\pi)\) and \(b=\ddm(\pi)\).  By the two-base
form, the normalized contribution of the \((a,b)\)-slice is bounded by
\[
  \Count_m(a,b)d^{m+3-a}S^{m+1-b}.
\]
Summing over the bidefect support and applying Theorem~\ref{thm:main} gives
\[
  4^{m+1}(d+\Loss_m)^{m+3}(S+\Loss_m)^{m+1}.
\]
Since \(k=m+1\), this is the displayed bound.
\end{proof}

\begin{remark}
The theorem is intentionally not stated as a full PPT threshold theorem.  The
analytic normalization, centering, and conversion from moment bounds to
spectral statements are external to the bidefect count.  The contribution of
this paper is the two-defect combinatorial estimate that the moment method
uses after the standard Wick expansion has supplied the termwise two-base
form.
\end{remark}

\section{Why Total Defect Is the Wrong Variable}

It is tempting to work only with
\[
  \dd(\pi)=\ddp(\pi)+\ddm(\pi).
\]
That variable is sufficient for the rank identity, but insufficient for the
scalar estimate.  If one only knows the total defect \(r\), the natural bound
has the form
\[
  \sum_{r}\#\{\dd=r\}D^{2m+4-r},
\]
where \(D\) must dominate both analytic bases.  In the exact square case this
may be harmless, but away from a common base it creates an artificial
max-base problem.  The bidefect decomposition avoids this: it preserves the
two pieces of information needed by the two binomial expansions
\[
  (d+\Loss_m)^{m+3},
  \qquad
  (S+\Loss_m)^{m+1}.
\]
The bookkeeping is therefore not cosmetic.  It is the combinatorial form in
which the analytic estimate is naturally true.

\bibliographystyle{alpha}
\bibliography{references.zotero}

\end{document}
